% Section 3: A* on the graph output.
\section{Solving the graph: A*}
\label{sec:astar}

\subsection{Why A* applies}

All doors weigh $1$, so Dijkstra's algorithm~\cite{dijkstra1959note} applies.
A*~\cite{hart1968formal} is Dijkstra with the queue ordered by
$f(v)=g(v)+h(v)$: $g$ the doors walked so far, $h$ an estimate of the doors
still to walk. With $h\equiv0$ it \emph{is} Dijkstra.

The estimate is the \keyterm{Manhattan distance} to the exit,
$h(x,y)=|x-x_e|+|y-y_e|$. Every step changes one coordinate by one, and walls
only remove edges, so $h$ never overestimates: it is admissible, and A* returns
a shortest route --- here, the unique one.

\subsection{How it solves the maze}

Each node holds its coordinates and its neighbours, which is all A* needs. It
settles cells in order of $f$, relaxing each neighbour through one door, and
stops when the exit is settled. Parent pointers give the route back.

\begin{algorithm}[H]
\caption{A* from $s$ to the exit $e$ on the graph $N$.}
\label{alg:astar}
\SetKwFunction{AStar}{AStar}
\SetKwFunction{Pop}{PopMin}
\SetKwFunction{Push}{Push}
\Fn{\AStar{$N$, $s$, $e$}}{
  $g[v]\leftarrow\infty$ for all $v$;\ $g[s]\leftarrow0$;\ $Q\leftarrow\{s\}$\;
  \While{$Q\neq\emptyset$}{
    $u\leftarrow$ \Pop{$Q$}\tcp*{least $g+h$}
    \lIf{$u=e$}{\textbf{break}}
    \ForEach{$v\in N.\textit{neighbours}(u)$}{
      \If{$g[u]+1<g[v]$}{
        $g[v]\leftarrow g[u]+1$;\ $\textit{parent}[v]\leftarrow u$;\ \Push{$Q$, $v$}\;
      }
    }
  }
  \Return the route from $e$ back along $\textit{parent}$ to $s$\;
}
\end{algorithm}

\subsection{Time complexity}

Two factors: how many cells are expanded, and what each queue operation costs.

\keyterm{Expansions: $\Theta(n)$.} In the worst case every cell is settled, and
in a maze that case is close to typical: $h$ prunes little, because a tree has
no rival routes to prune, only dead ends whose depth plus estimate already
exceeds the answer. Section~\ref{sec:comparison} measures $82\%$ of the cells.

\keyterm{Per operation: $O(1)$, not $O(\log n)$.} With a binary heap each
operation costs $\log F$, where $F$ is the queue size --- and with unit weights
the queue holds only the current frontier, measured at $60$ cells for
$n=65\,536$ ($F\approx n^{0.26}$). But the heap is unnecessary. The keys
$f=g+h$ are small integers and, $h$ being consistent, the popped $f$ never
decreases; a bucket queue with a moving pointer does each push and pop in
$O(1)$. With $h\equiv0$ this is breadth-first search with a FIFO.

\begin{keypoint}
A* on the maze runs in \keyterm{$\Theta(n)$}: linear in the cells, and in
practice touching most of them.
\end{keypoint}
